% CHƯƠNG 7 - Rubric: Phân tích mở rộng; Khả năng mở rộng, liên hệ thực tế (1,5 điểm)
\chapter{Cải tiến và mở rộng}
\label{chuong:caitien}

\section{Đề xuất cải tiến}

Hai cải tiến dưới đây nằm \textbf{ngoài phạm vi bài báo gốc}, do báo cáo này tự đề xuất và
hiện thực hóa, dựa trên quan sát rằng cả hai bộ dữ liệu đều dùng kiến trúc hồi quy đơn giản
nhất trong họ của nó (mục~\ref{sec:rnn}).

\subsection{Lý thuyết bổ sung: LSTM, GRU và LSTM hai chiều}
\label{sec:lythuyet-caitien}

Hai cải tiến của chương này cần thêm lý thuyết mà bài báo gốc không dùng, nên không đặt ở
Chương~\ref{chuong:nentang}. Long Short-Term Memory (LSTM) \cite{hochreiter1997lstm} khắc phục
nhược điểm tiêu biến gradient của RNN (mục~\ref{sec:rnn}) bằng cách thêm một trạng thái ô nhớ
(cell state) $c_t$ và ba cổng điều khiển luồng thông tin: cổng quên (forget gate) $f_t$, cổng
vào (input gate) $i_t$ và cổng ra (output gate) $o_t$:
\begin{align}
	f_t &= \sigma\!\left(W_f [h_{t-1}, x_t] + b_f\right), \\
	i_t &= \sigma\!\left(W_i [h_{t-1}, x_t] + b_i\right), \\
	o_t &= \sigma\!\left(W_o [h_{t-1}, x_t] + b_o\right), \\
	\tilde{c}_t &= \tanh\!\left(W_c [h_{t-1}, x_t] + b_c\right), \\
	c_t &= f_t \odot c_{t-1} + i_t \odot \tilde{c}_t, \\
	h_t &= o_t \odot \tanh(c_t),
	\label{eq:lstm}
\end{align}
trong đó $\sigma$ là hàm sigmoid, $\odot$ là tích theo từng phần tử, và $[h_{t-1}, x_t]$ là
phép nối hai véc-tơ. Cổng quên cho phép mô hình giữ lại hoặc loại bỏ thông tin trong ô nhớ một
cách có chọn lọc, giúp gradient lan truyền qua nhiều bước thời gian mà không bị tiêu biến.

Gated Recurrent Unit (GRU) \cite{cho2014gru} là một biến thể đơn giản hơn LSTM, gộp cổng quên
và cổng vào thành một cổng cập nhật (update gate) $z_t$, và bỏ hẳn ô nhớ riêng, chỉ dùng trạng
thái ẩn $h_t$:
\begin{align}
	z_t &= \sigma\!\left(W_z [h_{t-1}, x_t] + b_z\right), \\
	r_t &= \sigma\!\left(W_r [h_{t-1}, x_t] + b_r\right), \\
	\tilde{h}_t &= \tanh\!\left(W_h [r_t \odot h_{t-1}, x_t] + b_h\right), \\
	h_t &= (1 - z_t) \odot h_{t-1} + z_t \odot \tilde{h}_t,
	\label{eq:gru}
\end{align}
trong đó $r_t$ là cổng đặt lại (reset gate). Với ít tham số hơn LSTM, GRU thường huấn luyện
nhanh hơn, nhưng do thiếu ô nhớ riêng, khả năng kiểm soát luồng thông tin dài hạn thường kém
linh hoạt hơn LSTM.

LSTM hai chiều (Bidirectional LSTM) chạy hai mạng LSTM độc lập trên cùng một chuỗi đầu vào:
một theo chiều thuận ($t=1 \to T$) và một theo chiều ngược ($t=T \to 1$), rồi nối
(concatenate) hai trạng thái ẩn tại mỗi bước thời gian. Nhờ đó, biểu diễn tại mỗi bước thời
gian mang thông tin của \textit{cả ngữ cảnh trước và sau} nó trong cùng một cửa sổ dữ liệu,
hữu ích khi toàn bộ cửa sổ đã có sẵn tại thời điểm suy luận (như một đoạn tín hiệu điện não đã
thu xong), khác với bài toán dự báo tuần tự từng bước, nơi chỉ có thông tin quá khứ.

\subsection{Đổi kiến trúc hồi quy RNN sang LSTM/GRU trên Mackey--Glass}

Tác giả dùng RNN nguyên bản (\texttt{nn.RNN}, mục~\ref{sec:rnn}) cho cả Teacher, Baseline và
Student trên Mackey--Glass --- kiến trúc đơn giản nhất trong họ mạng hồi quy, dễ gặp hiện tượng
tiêu biến gradient hơn LSTM hoặc GRU trên chuỗi dài. Báo cáo đề xuất thay kiến trúc này bằng
LSTM hoặc GRU, áp dụng \textit{đồng thời} cho cả ba mô hình (giữ nguyên tắc ``một họ kiến trúc
dùng chung'' của bài báo, chỉ đổi loại cell hồi quy), nhằm kiểm tra: kiến trúc mạnh hơn có giúp
FGL cải thiện nhiều hơn không, hay mức cải thiện tương đối giữa Baseline và Student là bất biến
theo kiến trúc.

\subsection{LSTM hai chiều cho CNN-LSTM trên CHB-MIT}

Tương tự, mô hình \texttt{CNN\_LSTM\_Model} dùng cho cả Teacher và Student/Baseline trên
CHB-MIT chỉ dùng LSTM một chiều. Vì mỗi cửa sổ tín hiệu EEG đưa vào mô hình đã được thu thập
đầy đủ trước khi xử lý (không giống bài toán dự báo tuần tự từng bước, nơi chỉ có thông tin quá
khứ tại thời điểm dự báo), mỗi bước thời gian trong cửa sổ \textit{hoàn toàn có thể} được biểu
diễn bằng cả ngữ cảnh trước và sau nó, điều mà LSTM một chiều không khai thác được. Báo cáo đề
xuất thay LSTM một chiều bằng LSTM hai chiều (mục~\ref{sec:lythuyet-caitien}) trong
\texttt{CNN\_LSTM\_Model},
thêm tham số \texttt{bidirectional} (mặc định tắt, không đổi hành vi gốc khi không bật), áp dụng
cho cả Teacher và Student.

\section{Kết quả thực nghiệm cải tiến}

\subsection{LSTM/GRU trên Mackey--Glass: đã có kết quả đầy đủ}

Bảng~\ref{bang:mg-lstm} và Bảng~\ref{bang:mg-gru} trình bày kết quả sweep sáu seed (cùng seed
như Bảng~\ref{bang:mg-sweep}) với kiến trúc LSTM và GRU.

\begin{table}[H]{Kết quả sweep sáu seed trên Mackey--Glass --- kiến trúc LSTM}
	\label{bang:mg-lstm}
	\centering
	\begin{tabular}{cccc c}
		\toprule
		\textbf{Seed} & \textbf{Teacher MSE} & \textbf{Baseline MSE} & \textbf{Student MSE} & \textbf{Giảm MSE} \\
		\midrule
		0  & 4,1850 & 145,4412 & 92,5198 & 36,39\% \\
		1  & 4,6782 & 138,9365 & 89,1096 & 35,86\% \\
		2  & 4,6281 & 142,7424 & 94,6453 & 33,70\% \\
		3  & 4,1689 & 139,3032 & 92,2201 & 33,80\% \\
		4  & 4,7799 & 143,8094 & 98,7154 & 31,36\% \\
		42 & 4,4145 & 144,6108 & 97,3967 & 32,65\% \\
		\midrule
		\textbf{Trung bình} & 4,476 & 142,477 & 94,101 & \textbf{33,96\% (độ lệch chuẩn 1,74)} \\
		\bottomrule
	\end{tabular}
\end{table}

\begin{table}[H]{Kết quả sweep sáu seed trên Mackey--Glass --- kiến trúc GRU}
	\label{bang:mg-gru}
	\centering
	\begin{tabular}{cccc c}
		\toprule
		\textbf{Seed} & \textbf{Teacher MSE} & \textbf{Baseline MSE} & \textbf{Student MSE} & \textbf{Giảm MSE} \\
		\midrule
		0  & 4,4406 & 143,1596 & 98,4052 & 31,26\% \\
		1  & 4,5018 & 142,7224 & 98,0985 & 31,27\% \\
		2  & 4,2662 & 140,9470 & 94,1331 & 33,21\% \\
		3  & 4,2195 & 146,1191 & 97,3077 & 33,41\% \\
		4  & 4,4276 & 143,0985 & 87,6893 & 38,72\% \\
		42 & 4,7373 & 142,0975 & 98,7189 & 30,53\% \\
		\midrule
		\textbf{Trung bình} & 4,432 & 142,197 & 95,726 & \textbf{33,07\% (độ lệch chuẩn 2,74)} \\
		\bottomrule
	\end{tabular}
\end{table}

Bảng~\ref{bang:mg-sosanh-kientruc} tổng hợp so sánh ba kiến trúc.

\begin{table}[H]{So sánh mức giảm MSE trung bình giữa ba kiến trúc hồi quy (trung bình sáu seed)}
	\label{bang:mg-sosanh-kientruc}
	\centering
	\begin{tabular}{lcc}
		\toprule
		\textbf{Kiến trúc} & \textbf{Giảm MSE trung bình} & \textbf{So với RNN} \\
		\midrule
		RNN (gốc tác giả) & 26,42\% (độ lệch chuẩn 3,86) & --- \\
		LSTM & \textbf{33,96\% (độ lệch chuẩn 1,74)} & \textbf{+7,54 điểm \%} \\
		GRU & \textbf{33,07\% (độ lệch chuẩn 2,74)} & \textbf{+6,65 điểm \%} \\
		\bottomrule
	\end{tabular}
\end{table}

Cả LSTM và GRU đều cải thiện trên toàn bộ sáu seed, không có ngoại lệ. LSTM nhỉnh hơn GRU cả về
mức cải thiện trung bình lẫn độ ổn định (độ lệch chuẩn nhỏ hơn), phù hợp với trực giác rằng cổng
quên/nhớ riêng của LSTM kiểm soát luồng thông tin tốt hơn trên chuỗi có tính hỗn loạn cao như
Mackey--Glass. Đáng chú ý, MSE của \textit{Teacher} cũng giảm rõ rệt khi đổi kiến trúc (RNN
$\approx 5,86$ trung bình $\to$ LSTM $\approx 4,48 \to$ GRU $\approx 4,43$), cho thấy kiến trúc
mạnh hơn giúp cả nhiệm vụ 1-bước của teacher, không chỉ nhiệm vụ $H$-bước của student.

\subsection{Kiểm tra độ bền theo điểm cắt dữ liệu: time-series cross-validation}
\label{sec:mg-cv}

Mục~\ref{sec:mg-sweep-seed} đã chỉ ra một hạn chế: sweep sáu seed chỉ thay đổi khởi tạo trọng
số ngẫu nhiên, còn điểm chia train/val/test trên chuỗi Mackey--Glass luôn cố định (60\%/20\%/
20\% đầu tiên của chuỗi). Để kiểm tra riêng ảnh hưởng của \textit{vị trí chia dữ liệu}, báo cáo
bổ sung một thực nghiệm \textit{time-series cross-validation} kiểu rolling-origin (mở rộng cửa
sổ, tương đương \texttt{sklearn.TimeSeriesSplit}): không xáo trộn dữ liệu (tránh rò dữ liệu
tương lai), mà dùng ba đoạn đầu chuỗi ngày càng dài (70\%, 85\%, 100\% chuỗi gốc) làm ba
``fold''; trong mỗi fold, tỉ lệ 60/20/20 cũ vẫn được áp dụng, nên tập test của mỗi fold nằm ở
một đoạn càng về sau trên chuỗi thời gian. Thực nghiệm chạy cả ba kiến trúc (RNN, LSTM, GRU)
trên cùng seed 42, kết quả trong Bảng~\ref{bang:mg-cv}.

\begin{table}[H]{Kết quả time-series cross-validation (rolling-origin, 3 fold) trên Mackey--Glass}
	\label{bang:mg-cv}
	\centering
	\begin{tabular}{lcccc}
		\toprule
		\textbf{Fold (\% dữ liệu)} & \textbf{RNN} & \textbf{LSTM} & \textbf{GRU} \\
		\midrule
		Fold 1 (70\%) & 24,39\% & 32,82\% & 27,31\% \\
		Fold 2 (85\%) & 29,47\% & 32,72\% & 30,32\% \\
		Fold 3 (100\%) & 25,10\% & 32,65\% & 30,53\% \\
		\midrule
		\textbf{Trung bình} & 26,32\% (độ lệch chuẩn 2,25) & \textbf{32,73\% (độ lệch chuẩn 0,07)} & 29,38\% (độ lệch chuẩn 1,47) \\
		\bottomrule
	\end{tabular}
\end{table}

Fold 3 (100\% dữ liệu) dùng đúng cấu hình của ``lần chạy chuẩn'' seed 42
(Chương~\ref{chuong:taihien}), nên kết quả ở dòng này khớp \textit{chính xác từng số lẻ} với
Bảng~\ref{bang:mg-chuan}, \ref{bang:mg-lstm} và \ref{bang:mg-gru} (ví dụ RNN 25,10\%, LSTM
32,65\%, GRU 30,53\%) -- một cách kiểm chứng chéo cho thấy thực nghiệm cross-validation được
cài đặt đúng, không có lỗi âm thầm làm lệch số liệu.

Kết quả quan trọng nhất: \textbf{cả ba kiến trúc đều cải thiện dương trên cả ba fold}, không
có trường hợp nào FGL làm tệ hơn baseline khi đổi điểm chia dữ liệu -- khác với thực nghiệm
quét horizon ở Bảng~\ref{bang:mg-drift}, nơi horizon=2 là ngoại lệ âm. Đáng chú ý hơn,
\textbf{LSTM có độ lệch chuẩn rất nhỏ giữa các fold (0,07 điểm \%)}, nhỏ hơn nhiều so với RNN
(2,25) và GRU (1,47) -- cho thấy mức cải thiện của LSTM không chỉ ổn định qua các seed khởi tạo
khác nhau (mục~\ref{sec:mg-sweep-seed}) mà còn ổn định qua các vị trí chia dữ liệu khác nhau
trên chuỗi, một bằng chứng bổ sung cho việc chọn LSTM làm cải tiến kiến trúc.

\subsection{LSTM hai chiều trên CHB-MIT: đã cài đặt, chưa có kết quả số}

Khác với cải tiến trên Mackey--Glass, cải tiến LSTM hai chiều cho CHB-MIT \textbf{mới hoàn
thành phần cài đặt mã nguồn} (tham số \texttt{bidirectional}, cờ \texttt{--bidirectional}, đã
kiểm chứng cơ chế LSTM và tầng kết nối đầy đủ phía sau hoạt động đúng với dữ liệu giả lập ở cả
hai chế độ), nhưng \textbf{chưa có kết quả AUC-ROC thực tế} để báo cáo, vì quá trình chạy lại
Giai đoạn 1 và Giai đoạn 3 trên Kaggle (Chương~\ref{chuong:taihien}) còn đang được hoàn thiện
tại thời điểm viết báo cáo này. Đây được nêu rõ là một phần \textbf{đang thực hiện}, không phải
kết quả đã có nhưng bị bỏ sót.

\section{Liên hệ thực tế}

Đánh đổi giữa tính hữu dụng và độ chính xác của dự báo sớm/muộn (Chương~\ref{chuong:baitoan})
không chỉ xuất hiện trong dự báo cơn động kinh, mà còn phổ biến ở nhiều lĩnh vực khác có dữ
liệu dạng chuỗi thời gian với khả năng thay đổi cấu trúc hoặc có biến cố cần cảnh báo sớm:

\begin{itemize}
	\item \textbf{Bảo trì dự đoán (predictive maintenance)} trong công nghiệp: tín hiệu cảm
	biến của máy móc thường có dấu hiệu bất thường nhẹ rất lâu trước khi hỏng hóc thực sự xảy
	ra; một mô hình kiểu FGL có thể dùng dữ liệu ngay-trước-khi-hỏng (dễ nhận biết, như teacher)
	để dẫn dắt một mô hình dự báo sớm hơn (khó nhận biết hơn, như student).
	\item \textbf{Dự báo tài chính có thay đổi chế độ (regime change)}: các chuỗi giá hoặc chỉ
	số tài chính thường có các giai đoạn với động lực khác nhau (như ví dụ về gãy cấu trúc do
	đại dịch COVID-19 trong nhiều nghiên cứu dự báo chuỗi thời gian); cơ chế teacher/student
	lệch cửa sổ thời gian tương tự Mackey--Glass (mục~\ref{sec:cochetuonglai}) có thể áp dụng
	trực tiếp cho bài toán này.
	\item \textbf{Cảnh báo sớm trong y tế} cho các biến cố khác ngoài động kinh, như nhiễm
	khuẩn huyết (sepsis) hay rối loạn nhịp tim, nơi cùng tồn tại đánh đổi giữa phát hiện sớm
	(hữu dụng, khó) và phát hiện muộn (dễ, ít giá trị can thiệp) --- cấu trúc bài toán giống
	hoàn toàn với dự báo cơn động kinh trên CHB-MIT.
\end{itemize}

Hai cải tiến đề xuất trong chương này --- đổi kiến trúc hồi quy mạnh hơn, và khai thác ngữ
cảnh hai chiều khi dữ liệu đầu vào đã có sẵn đầy đủ --- đều là các thay đổi kiến trúc đơn giản,
không đụng đến logic chưng cất tri thức cốt lõi của FGL, nên có thể kỳ vọng mang lại lợi ích
tương tự khi áp dụng FGL cho các lĩnh vực khác kể trên.
